\section{Reproducibility statement}\label{sec:repro}
% source: RM:14, 174, 228, 230; TDR:267-274; RP:190-196; EXP:939-968; XR:99-132; KCL:46-52; AUD:61-70; HER:57-65; L1:62-68; L2:67-70.
Anyone can run the training and evaluation harness and its 1,638 collected tests without a GPU, evaluate the pipeline end to end on the public development fixtures with an echo backend, and check every configuration hash in this paper except the pilot's against the committed configurations. Anyone can rebuild the data: at the recorded commit of the companion repository, the public pipeline reproduces release v0.0.6 byte for byte (content hash \code{3cc9a744\ldots}) and, from the reasoning change on, release v0.0.7 (\code{b53afa42\ldots}); the test split and the benchmark then carry the SHA-256s cited in Section~\ref{sec:data}. Holders of a release can retrain and re-evaluate each adapter from the runbooks in the repository with the recorded configurations, commits and seeds, and compare their adapters against the recorded SHA-256s and their results against the frozen result-file hashes (Appendix~\ref{app:artifacts}). Nobody can reproduce the pilot's base revision, which was never recorded (deviation D7); its run is reproducible only if the upstream \code{main} pointer has not moved. Every run's manifest records its configuration hash, git commit, dataset hash and runtime adjustments (none), and each verdict in this paper is traceable to a dated entry in the research record. The repository's history was rewritten before publication to remove assistant co-author trailers from seven commit messages; the content, authors and dates of the commits are unchanged, and a table in the record maps each identifier cited by the frozen records to its public counterpart, six of which appear in Appendix~\ref{app:artifacts}.
